\BourbakiExerciseGroup

All fields in this section are commutative unless explicitly stated otherwise.

\begin{enumerate}
\def\labelenumi{\arabic{enumi})}
\item
  Let u be an endomorphism of a finite dimensional vector space E, and let V be a subspace of E invariant under u. Show that the characteristic polynomial of the restriction of u to V divides that of u. Hence obtain an elementary proof of the Cayley-Hamilton theorem (reduce to the case where \(E_{u}\) is cyclic, and use formula (2) of VII, p.~30).
\item
  \begin{enumerate}
  \def\labelenumii{\alph{enumii})}
  \tightlist
  \item
    Show that every square matrix over a field K is similar to its transpose (reduce to the case of a Jordan matrix).
  \end{enumerate}
\end{enumerate}

\begin{enumerate}
\def\labelenumi{\alph{enumi})}
\setcounter{enumi}{1}
\item
  Show that over the ring \(\mathbf{Z}\) of rational integers, the matrix \(\begin{pmatrix} 8 & 2 \\ 0 & 1 \end{pmatrix}\) is not similar to its transpose.
\item
  Let \(u\) be an endomorphism of a finite dimensional vector space over \(\mathbf{K}\) . Then every eigenvalue \(\lambda\) of \(u\) is also an eigenvalue of the transpose endomorphism \(^t u\) of \(E^*\) . Moreover, if \(\lambda\) and \(\mu\) are two distinct eigenvalues of \(u\) , if \(x\) is an eigenvector of \(u\) corresponding to \(\lambda\) , and if \(y^*\) is an eigenvector of \(^t u\) corresponding to \(\mu\) , then \(\langle x, y^* \rangle = 0\) .
\end{enumerate}

\begin{enumerate}
\def\labelenumi{\arabic{enumi})}
\setcounter{enumi}{2}
\item
  Let u be an endomorphism of a finite dimensional vector space E, let \(\lambda\) be an eigenvalue of u, and let v be the endomorphism \(u - \lambda\) . Show that E is the direct sum of \(v^{-1}(0)\) and \(v(E)\) if and only if \(\lambda\) is a simple root of the minimal polynomial of u.
\item
  Show, by adjoining a transcendental element to the base field, that formula (6) in Prop. 10 of VII, p.~36 can be deduced from formula (8). Prove the latter directly by decomposing the polynomial q into linear factors.
\item
  \begin{enumerate}
  \def\labelenumii{\alph{enumii})}
  \tightlist
  \item
    Let K be an algebraically closed field, and let A be a square matrix of order n over K, in the form of a diagonal block \(\text{diag}(A_i)\) of lower triangular matrices \(A_i\) , such that the diagonal entries of \(A_i\) are all equal to the same element \(\alpha_i\) of K ( \(1 \leqslant i \leqslant r\) ), and \(\alpha_i \neq \alpha_j\) for \(i \neq j\) . Show that every square matrix B which commutes with A is a diagonal block \(\text{diag}(B_i)\) , where \(B_i\) is the same size as \(A_i\) for \(1 \leqslant i \leqslant r\) .
  \end{enumerate}
\end{enumerate}

\begin{enumerate}
\def\labelenumi{\alph{enumi})}
\setcounter{enumi}{1}
\tightlist
\item
  Let M be a set of commuting square matrices of order n over K. Show that by applying the same similarity to all the matrices in M, we can make them all into diagonal blocks \(\text{diag}(X_i)\) of lower triangular matrices, where the number r of \(X_i\) 's and the order \(m_i\) of \(X_i\) for each i, are the same for each matrix in M, and where all the eigenvalues of \(X_i\) are equal. (Use a) to reduce to the case where r = 1, then Lemma 3 of VII, p.~41 to show that in this case all the matrices in M have nonzero common eigenvector, and complete the proof by induction on n.)
\end{enumerate}

\begin{enumerate}
\def\labelenumi{\arabic{enumi})}
\setcounter{enumi}{5}
\item
  Two pairs \((A_{1}, A_{2})\) and \((B_{1}, B_{2})\) of square matrices of order n over a commutative field K are said to be equivalent if there exist two invertible square matrices P and Q over K such that \(A_{1} = PB_{1}Q\) and \(A_{2} = PB_{2}Q\) . If \(A_{1}\) and \(B_{1}\) are invertible, show that the pairs are equivalent if and only if the matrices \(XA_{1} + A_{2}\) and \(XB_{1} + B_{2}\) are equivalent over the ring K{[}X{]}.
\item
  Let \(A\) be a matrix with \(m\) rows and \(n\) columns, of rank \(r\) , over a field \(K\) . If \(P\) is a square matrix of order \(m\) and \(Q\) is a square matrix of order \(n\) over \(K\) , such that \(PAQ = A\) , show that there exists a matrix \(P'\) similar to \(P\) and a matrix \(Q'\) similar to \(Q\) such that
\end{enumerate}

\[
P ^ {\prime} = \left( \begin{array}{c c} R & S \\ 0 & U \end{array} \right), \quad Q ^ {\prime} = \left( \begin{array}{c c} R ^ {- 1} & 0 \\ S ^ {\prime} & U ^ {\prime} \end{array} \right)
\]

where R is an invertible square matrix of order r, U (resp. \(U'\) ) is a square matrix of order m - r (resp. n - r), and S (resp. \(S'\) ) is a matrix with r rows and m - r columns (resp. n - r rows and m columns).

\begin{enumerate}
\def\labelenumi{\arabic{enumi})}
\setcounter{enumi}{7}
\item
  Let E be a vector space of finite dimension n over a field K, let u be an endomorphism of E, let V be a subspace of E invariant under u, and let \(u|V\) be the restriction of u to V. If \(f_{i}\) ( \(1 \leqslant i \leqslant n\) ) (resp. \(g_{j}\) ( \(1 \leqslant j \leqslant m\) )) are the similarity invariants of u (resp. \(u|V\) ), show that \(g_{j}\) divides \(f_{j+n-m}\) for \(1 \leqslant j \leqslant m\) (cf.~VII, p.~64, Ex. 8, d)). State and prove the analogous property of the endomorphism of E/V induced by u.
\item
  Let \(A\) be a square matrix of order \(n\) , let \(B\) be a square matrix or order \(m\) , and let \(C\) be a matrix with \(m\) rows and \(n\) columns over a field \(K\) . Reduce the equation \(X A = B X + C\) , where \(X\) is an unknown matrix with \(m\) rows and \(n\) columns over \(K\) , to a system of linear congruences in the polynomial ring \(K[Y]\) (for all \(f \in K[Y]\) , notice that \(X f(A) = f(B)X + H(f)\) for some well defined matrix \(H(f)\) with \(m\) rows and \(n\) columns; let \(\oplus_{i} E_i\) be the decomposition of the module \(E\) associated to \(A\) (Sect. 1) as a direct sum of
\end{enumerate}

cyclic modules, whose annihilators are the similarity invariants \(f_{i}\) of A which are distinct from 1 (Prop. 2), and let \(a_{i}\) be generators of the \(E_{i}\) ; in the same way define the submodules \(F_{j}\) , generated by \(b_{j}\) , and the polynomials \(g_{j}\) corresponding to B; write \(Xf_{i}(A)a_{i}=0\) ; then decompose \(Xa_{i}\) and \(H(f_{i})a_{i}\) in terms of the \(F_{j}\) , and show that the conditions thus obtained from the equations \(Xf_{i}(A)a_{i}=0\) are necessary and sufficient for the existence of a solution).

\begin{enumerate}
\def\labelenumi{\arabic{enumi})}
\setcounter{enumi}{9}
\tightlist
\item
  \begin{enumerate}
  \def\labelenumii{\alph{enumii})}
  \tightlist
  \item
    Show that, over an arbitrary field K, the matrix
  \end{enumerate}
\end{enumerate}

\[
\left( \begin{array}{c c c c c c} \lambda & \alpha & \beta_ {13} & \beta_ {14} & ... & \beta_ {1n} \\ 0 & \lambda & \alpha & \beta_ {24} & ... & \beta_ {2n} \\ 0 & 0 & \lambda & \alpha & ... & \beta_ {3n} \\ \cdots & \cdots & \cdots & \cdots & \cdots & \cdots \\ 0 & 0 & 0 & 0 & ... & \alpha \\ 0 & 0 & 0 & 0 & ... & \lambda \end{array} \right)
\]

is similar to the Jordan matrix \(U_{n,\lambda}\) if \(\alpha \neq 0\) .

\begin{enumerate}
\def\labelenumi{\alph{enumi})}
\setcounter{enumi}{1}
\item
  If \(\lambda \neq 0\) and \(\mathbf{K}\) is of characteristic 0, show that \((U_{n,\lambda})^m\) is similar to the Jordan matrix \(U_{n,\lambda^m}\) for every integer \(m\) (positive or negative).
\item
  Show that \((U_{n,0})^m = 0\) for \(m \geqslant n\) . For \(0 < m < n\) , let \(n - 1 = km + q\) with \(0 \leqslant q < m\) ; show that \((U_{n,0})^m\) has \(q + 1\) similarity invariants equal to \(\mathbf{X}^{k+1}\) and \(m - q - 1\) similarity invariants equal to \(\mathbf{X}^k\) (arrange the canonical basis of \(\mathbf{K}^n\) in a suitable order).
\item
  Similarly, determine the similarity invariants of \(f(U_{n,\lambda})\) for \(f \in \mathbf{K}[\mathbf{X}]\) .
\item
  Discuss the same problems for K of characteristic \(p \neq 0\) .
\end{enumerate}

\begin{enumerate}
\def\labelenumi{\arabic{enumi})}
\setcounter{enumi}{10}
\item
  Let K be an algebraically closed field, let A be an invertible matrix over K, and m \textgreater{} 0 an integer. Show that, if m is not a multiple of the characteristic of K, then there exists a polynomial \(g(\mathbf{X}) \in \mathbf{K}[\mathbf{X}]\) such that \((g(A))^{m} = A\) (use the Cayley-Hamilton theorem and Ex. 28 of VII, p.~54). Generalise to the equation \(f(U) = A\) ( \(f \in K[X]\) , and U an unknown \(n \times n\) matrix over K).
\item
  \begin{enumerate}
  \def\labelenumii{\alph{enumii})}
  \tightlist
  \item
    Let \(A\) and \(B\) be two square matrices of order \(n\) over a commutative field \(\mathbf{K}\) , and let \(f_i\) and \(g_i\) ( \(1 \leqslant i \leqslant n\) ) be their respective similarity invariants. Show that the vector space (over \(\mathbf{K}\) ) of square matrices \(U\) , of order \(n\) , such that \(UA = BU\) , is isomorphic to the space \(\mathbf{M}_0 = \mathbf{M} / \mathbf{N}\) defined as follows: \(\mathbf{M}\) is the space of square matrices \((u_{ij}(\mathbf{X}))\) of order \(n\) over \(\mathbf{K}[\mathbf{X}]\) such that \(u_{ij}(\mathbf{X}) f_j(\mathbf{X})\) is a multiple of \(g_i(\mathbf{X})\) for each pair \(i, j\) of indices; and \(\mathbf{N}\) is the subspace of \(\mathbf{M}\) consisting of those matrices for which \(u_{ij}(\mathbf{X})\) is a multiple of \(g_i(\mathbf{X})\) for each pair \(i, j\) of indices. (Consider \(U, A\) and \(B\) as the matrices of endomorphisms \(u, v\) and \(w\) respectively of \(E = K^n\) , and notice that \(u\) is a \(K[X]\) -linear map from \(E_v\) into \(E_w\) ; now express \(E_v\) and \(E_w\) as quotient modules.)
  \end{enumerate}
\end{enumerate}

\begin{enumerate}
\def\labelenumi{\alph{enumi})}
\setcounter{enumi}{1}
\item
  When B = A, the matrices U such that UA = AU form a ring P; then M is a ring of square matrices, and N a 2-sided ideal of M, such that P is isomorphic to the quotient ring M/N. Show that P is identical to the ring K{[}A{]} generated by \(I_{n}\) and A if and only if the minimal polynomial of A is equal to its characteristic polynomial.
\item
  For arbitrary \(A\) and \(B\) , show that the dimension of \(\mathbf{M}_0\) is equal to \(\sum_{i,j} m_{ij}\) , where
\end{enumerate}

\(m_{ij}\) is the degree of the gcd of \(f_i\) and \(g_j\) (same method as \(a\) ).

\begin{enumerate}
\def\labelenumi{\alph{enumi})}
\setcounter{enumi}{3}
\tightlist
\item
  Show that the greatest of the ranks of the matrices \(U \in \mathbf{M}_0\) is equal to \(\sum_{k} d_k\) , where
\end{enumerate}

\(d_{k}\) is the degree of the gcd of \(f_{k}\) and \(\pmb{g}_{k}\) (use Ex. 8).

\begin{enumerate}
\def\labelenumi{\alph{enumi})}
\setcounter{enumi}{4}
\tightlist
\item
  Generalise the results of a), c) and d) to the case where A is a square matrix of order n, B is a square matrix of order m, and U is a matrix with m rows and n columns.
\end{enumerate}

\begin{enumerate}
\def\labelenumi{\arabic{enumi})}
\setcounter{enumi}{12}
\tightlist
\item
  Let \((\mathbf{Z}_{ij})\) ( \(1 \leqslant i, j \leqslant n\) ) be \(n^2\) indeterminates over a field \(\mathbf{K}_0\) . Let \(\mathbf{A}\) be the polynomial ring \(\mathbf{K}_0[\mathbf{Z}_{11}, \ldots, \mathbf{Z}_{nn}]\) , and let \(\mathbf{K} = \mathbf{K}_0(\mathbf{Z}_{ij})\) be its field of fractions; let \(\chi_U\) be the characteristic polynomial of the matrix \(U = (\mathbf{Z}_{ij})\) over \(\mathbf{K}\) .
\end{enumerate}

\begin{enumerate}
\def\labelenumi{\alph{enumi})}
\item
  Show that \(\chi_{U}\) is a separable irreducible polynomial (show that, if \(\chi_{U}\) were reducible over K, it would be the product of two polynomials in \(K_{0}[X, Z_{ij}]\) of degree \textgreater0 in X; hence deduce that the characteristic polynomial of every square matrix of order n over K would be reducible; then derive a contradiction from the example of a polynomial with algebraically independent coefficients over \(K_{0}\) ; use the analogous method to prove separability).
\item
  Show that U is similar to a matrix in which every entry is equal to 0, 1 or a coefficient of the polynomial \(\chi_{U}\) (cf.~VII, p.~30, formula (2)). Deduce that, if \(K_{0}\) is an infinite field, and p is a polynomial in A such that \(p(u_{ij}^{\prime}) = p(u_{ij}^{\prime\prime})\) whenever the matrices \((u_{ij}^{\prime})\) and \((u_{ij}^{\prime\prime})\) over \(K_{0}\) are similar, then p belongs to the subring of A generated by the coefficients of \(\chi_{U}\) . State and prove the analogous result for rational functions \(r \in K = K_{0}(Z_{ij})\) .
\end{enumerate}

\begin{enumerate}
\def\labelenumi{\arabic{enumi})}
\setcounter{enumi}{13}
\tightlist
\item
  Let K be an imperfect field of characteristic p, let a be an element of \(K - K^{p}\) , and let L be the field \(\mathrm{K}(a^{1/p})\) . Consider the K-vector space \(E = L \otimes_{K} L\) ; let u (resp. \(u'\) ) be the K-endomorphism of E given by multiplication by \(a^{1/p} \otimes 1\) (resp. \(1 \otimes a^{1/p}\) ) in the algebra E. Show that u and \(u'\) are semi-simple, that u and \(u'\) commute, and that \(u - u'\) is nonzero nilpotent.
\end{enumerate}

